% TOP_PROBLEM: {"id": 2593, "title": "Kourovka Notebook Problem 21.84", "category_id": 17, "category": {"id": 17, "name": "group_theory", "display_name": "Group Theory", "description": "Problems about groups, group actions, representations, and related algebraic structures.", "slug": "group-theory", "order_index": 17, "created_at": "2026-07-31T15:26:25.671Z"}, "status": "open", "problem_number": "KOU-21.84", "set_id": 10, "statusQualification": "Distinguishes matrix size d from permutation degree n and corrects the reversed metric arguments in the source transcription."}
% FIRST_POSTED: 2026-10-01
% ENTRY_KIND: statement_only
% UnsolvedMath ID 2593; source catalogue code KOU-21.84
% !TeX program = lualatex
% Definition and sources only; this file is not an LLM solution attempt.
\documentclass[11pt,a4paper]{article}
\usepackage[margin=25mm]{geometry}
\usepackage{fontspec}
\setmainfont{lmroman10-regular.otf}[BoldFont=lmroman10-bold.otf,
 ItalicFont=lmroman10-italic.otf,BoldItalicFont=lmroman10-bolditalic.otf]
\usepackage{amsmath,amssymb,mathtools}
\usepackage{unicode-math}
\setmathfont{latinmodern-math.otf}
\AtBeginDocument{\let\setminus\smallsetminus}
\usepackage{xurl,hyperref}
\hypersetup{colorlinks=true,urlcolor=blue}
\setcounter{secnumdepth}{0}
\setlength{\parindent}{0pt}
\setlength{\parskip}{5pt}
\setlength{\emergencystretch}{3em}
\begin{document}
\section*{Kourovka Notebook Problem 21.84}\label{2593}
{\small UnsolvedMath ID 2593 · KOU-21.84 · Group Theory · Kourovka Notebook - New Problems, Issue 21}
\subsection{Definitions and mathematical statement}
For an integer $d\ge1$, let $G=\operatorname{SL}_d(\mathbb Z)$ be the group of $d$ by $d$ integer matrices of determinant one. The matrix size $d$ is fixed while the permutation degree below tends to infinity. Let $S_n$ act on $[n]=\{1,\ldots,n\}$, with normalized Hamming distance
\[
 d_n(\sigma,\tau)=n^{-1}|\{i\in[n]:\sigma(i)\ne\tau(i)\}|.
\]
An almost-homomorphism is a sequence of maps $f_n:G\to S_n$ such that
\[
 d_n(f_n(g)f_n(h),f_n(gh))\longrightarrow0
 \quad\text{for every fixed }g,h\in G.
\]

For $m\ge n$, $\sigma\in S_n$, and $\tau\in S_m$, identify $[n]\subseteq[m]$ and put
\[
 d_n^{\mathrm{flex}}(\sigma,\tau)
 =n^{-1}\bigl(|\{i\in[n]:\sigma(i)\ne\tau(i)\}|+m-n\bigr).
\]
Is $\operatorname{SL}_d(\mathbb Z)$ flexibly permutation-stable? Explicitly, does every almost-homomorphism admit integers $m_n\ge n$ and genuine homomorphisms $\rho_n:G\to S_{m_n}$ such that
\[
 d_n^{\mathrm{flex}}(f_n(g),\rho_n(g))\longrightarrow0
 \qquad\text{for every fixed }g\in G?
\]
The problem asks this for the matrix groups as $d$ varies; the case $d=1$ is the trivial boundary case. Flexible convergence requires both approximation on the original points and $(m_n-n)/n\to0$. Thus the correction may enlarge the underlying sets by a vanishing relative amount, while satisfying the group law exactly. No uniformity over all group elements is imposed on these pointwise limits.
\subsection{Short English statement}
Can every approximate finite permutation action of $\operatorname{SL}_d(\mathbb Z)$ be corrected after adding only a vanishing proportion of extra points?
\subsection{Sources}
\begin{itemize}
\item[S1] ULAM.AI and the UnsolvedMath Contributors, supplied problems.json, record 2593 (KOU-21.84), Kourovka Notebook - New Problems, Issue 21. Catalogue identity and source transcription; CC BY 4.0.
\url{https://huggingface.co/datasets/ulamai/UnsolvedMath}.
\item[S2] The Kourovka Notebook, 21st edition, new problems, Problem 21.84. Expanded formulation of the supplied catalogue statement.
\url{https://arxiv.org/pdf/1401.0300v47}.
\end{itemize}
\end{document}
